\chapter{Proyecto integrador: LibreReserva operable}
\label{chap:proyecto}

\begin{objetivos}
El estudiante integrará arquitectura, código, entrega, operación, seguridad e IA responsable en un producto pequeño, defendible y reproducible. El objetivo no es acumular tecnologías, sino demostrar decisiones proporcionadas mediante evidencia.
\end{objetivos}

\section{Problema y restricciones}

LibreReserva permite publicar recursos institucionales, consultar disponibilidad, reservar, cancelar y auditar cambios. El equipo debe construir un incremento operable para dos perfiles: usuario y administrador. Se asumen datos sintéticos, presupuesto de infraestructura nulo y ejecución local con software libre.

Restricciones iniciales: API HTTP documentada; persistencia PostgreSQL; ejecución en contenedor OCI sin privilegio; automatización en Forgejo; despliegue local reproducible; telemetría OpenTelemetry; modelado de amenazas; registro de uso de IA. El equipo puede sustituir una herramienta por otra libre si mantiene la capacidad y justifica la decisión en un ADR.

\begin{figure}[htbp]
\centering
\begin{tikzpicture}[node distance=1.0cm and 1.2cm, every node/.style={font=\small}]
\node[actor] (u) {Usuario};
\node[system, right=of u] (api) {API modular};
\node[evidence, right=of api] (db) {PostgreSQL};
\node[system, below=of api] (worker) {Procesador de eventos};
\node[evidence, right=of worker] (mq) {RabbitMQ};
\node[activity, above=of api] (ci) {Forgejo Actions};
\node[evidence, above=of db] (obs) {Telemetría};
\draw[arrow] (u)--(api); \draw[arrow] (api)--(db); \draw[arrow] (api)--(mq);
\draw[arrow] (mq)--(worker); \draw[arrow] (ci)--(api); \draw[arrow] (api)--(obs);
\end{tikzpicture}
\caption{Arquitectura de referencia, deliberadamente sustituible mediante ADR.}
\end{figure}

\section{Resultados de aprendizaje y evidencias}

\begin{longtable}{p{1.2cm}p{5.1cm}p{8cm}}
\caption{Matriz de trazabilidad del proyecto.}\\
\toprule
\textbf{RA} & \textbf{Capacidad} & \textbf{Evidencia obligatoria}\\
\midrule
\endfirsthead
\toprule
\textbf{RA} & \textbf{Capacidad} & \textbf{Evidencia obligatoria}\\
\midrule
\endhead
RA1 & Analizar problema y restricciones & Mapa de actores, alcance, supuestos y riesgos.\\
RA2 & Priorizar atributos de calidad & Tres escenarios medibles y tabla de compromisos.\\
RA3 & Comunicar arquitectura & C4 de contexto/contenedores, UML selectivo y ADR.\\
RA4 & Diseñar modularidad & Límites, puertos, dependencias y prueba de arquitectura.\\
RA5 & Integrar datos y servicios & Contrato API, migración, idempotencia y prueba de integración.\\
RA6 & Construir y validar & Incremento vertical, suite y reporte de calidad.\\
RA7 & Entregar y operar & Canalización, imagen, manifiestos y despliegue reproducible.\\
RA8 & Observar confiabilidad & SLI/SLO, panel, trazas y experimento de fallo.\\
RA9 & Diseñar seguridad & Modelo de amenazas, SBOM, escaneos y triage.\\
RA10 & Usar IA responsablemente & Registro, evaluación comparada, revisión y defensa.\\
\bottomrule
\end{longtable}

\section{Iteraciones}

\begin{enumerate}
  \item \textbf{Descubrimiento y arquitectura.} Problema, restricciones, escenarios de calidad, C4, modelo de dominio y dos ADR. Validación: revisión cruzada y prototipo de riesgo.
  \item \textbf{Incremento ejecutable.} Crear y consultar reservas con persistencia, migración, manejo de conflictos y pruebas. Validación: contrato, cobertura razonada y demostración desde clon limpio.
  \item \textbf{Entrega segura.} Forgejo Actions, imagen OCI, SBOM, secretos, escaneos y despliegue kind. Validación: digest trazable y puerta que rechaza un cambio deliberadamente defectuoso.
  \item \textbf{Operación y resiliencia.} Telemetría, SLO, panel, alerta y experimento controlado. Validación: diagnóstico con métrica, traza y registro.
  \item \textbf{Evaluación de IA y cierre.} Experimento asistido, deuda, costo, riesgos pendientes, manual reproducible y defensa individual.
\end{enumerate}

Cada iteración termina con una versión demostrable. La retroalimentación modifica el producto y, cuando corresponde, el ADR o el modelo de amenazas. No se aceptan diagramas que describan un sistema distinto del entregado.

\section{Definición de terminado}

Una historia está terminada cuando: criterios de aceptación aprobados; revisión de otra persona; pruebas relevantes automatizadas; contrato y documentación actualizados; telemetría suficiente; amenazas examinadas; ningún secreto confirmado; canalización verde; imagen identificada por digest; y autoría/asistencia de IA declarada. Una excepción debe registrarse como deuda con riesgo, responsable y fecha.

\section{Rúbrica}

\begin{table}[htbp]
\centering
\caption{Rúbrica analítica del proyecto final.}
\small
\begin{tabularx}{\textwidth}{p{4cm}cX}
\toprule
\textbf{Dimensión} & \textbf{Peso} & \textbf{Desempeño esperado}\\
\midrule
Problema y decisiones arquitectónicas & 20\% & Restricciones explícitas, alternativas reales, escenarios medibles y ADR coherentes.\\
Construcción y pruebas & 20\% & Incremento correcto, modular, legible y probado por riesgo; datos reproducibles.\\
Entrega y operación & 15\% & Canalización confiable, artefacto trazable, despliegue y recuperación demostrados.\\
Observabilidad y confiabilidad & 15\% & SLI/SLO defendible, señales correlacionadas y experimento que sustenta una conclusión.\\
Seguridad, privacidad y suministro & 15\% & Amenazas priorizadas, controles probados, secretos protegidos, SBOM y triage.\\
IA responsable y trazabilidad & 10\% & Uso declarado, contexto protegido, evaluación independiente y comprensión demostrada.\\
Comunicación y reproducibilidad & 5\% & Documentación concisa, fuentes, licencias y ejecución desde entorno limpio.\\
\bottomrule
\end{tabularx}
\end{table}

Una falla crítica de integridad ---credenciales reales, datos personales sin autorización, resultados fabricados o atribución deliberadamente falsa--- se gestiona según el reglamento institucional y no se compensa con otros porcentajes.

\section{Validación final}

\begin{validacion}
Un evaluador parte de un clon limpio y sigue el README. Debe poder: crear entorno; ejecutar pruebas; construir imagen; generar SBOM; desplegar; realizar una reserva; observar su traza; ejecutar un caso adverso autorizado y relacionar la versión desplegada con una confirmación. Todo comando dependiente del sistema indica plataforma y versión.
\end{validacion}

\section{Defensa}

La demostración dura quince minutos: problema y decisión principal; recorrido funcional; fallo inducido y diagnóstico; control de seguridad; resultado del experimento con IA; deuda reconocida. Luego cada integrante responde sobre una parte seleccionada al azar y realiza un cambio pequeño. Esto evalúa comprensión compartida y reduce la utilidad de entregas que el equipo no puede explicar.

\begin{ejercicio}
\textbf{Entrega final.} Repositorio, etiqueta de versión, digest de imagen, paquete de evidencias y breve informe arquitectónico. Incluya licencia del proyecto, atribuciones de componentes, instrucciones de instalación, matriz RA--evidencia y registro de decisiones. No entregue dependencias descargadas, modelos de IA ni secretos dentro del repositorio.
\end{ejercicio}

\section{Cierre}

La arquitectura moderna es una práctica de decisiones comprobables. Las plataformas, incluida la IA generativa, cambian con rapidez; la capacidad durable consiste en formular restricciones, comparar alternativas, construir ciclos de evidencia y responder responsablemente por el sistema resultante.
